\documentclass[12pt]{article}
\usepackage[utf8]{inputenc}
\usepackage[T1]{fontenc}
\usepackage[spanish,mexico]{babel}
\usepackage{amsmath,amssymb,amsthm,mathtools}
\usepackage{geometry}
\usepackage{enumitem}
\geometry{letterpaper,margin=2.5cm}
\setlength{\parindent}{0pt}
\setlength{\parskip}{0.55em}
\newtheorem{teorema}{Teorema}[section]
\newtheorem{proposicion}[teorema]{Proposición}
\theoremstyle{definition}
\newtheorem{ejemplo}{Ejemplo}[section]
\newtheorem{ejercicio}{Ejercicio}[section]
\theoremstyle{remark}
\newtheorem{observacion}{Observación}[section]
\newcommand{\C}{\mathbb{C}}

\title{\textbf{Variable Compleja}\\[0.2cm]\Large Raíces de polinomios y Teorema Fundamental del Álgebra}
\author{Carlos Ernesto Martínez Rodríguez}
\date{}

\begin{document}
\maketitle
\tableofcontents
\newpage

\section{Propósito}
Estas notas complementarias fortalecen el cálculo de raíces de ecuaciones
polinomiales en $\C$ y utilizan los ejemplos para ilustrar el Teorema
Fundamental del Álgebra (TFA). La dificultad es deliberadamente mayor que en
los casos elementales $z^n=1$, $z^n=-1$ y $z^n=r$, pero se conservan
estructuras que permiten una solución completa en clase.

\section{Recordatorio: raíces $n$-ésimas}
Si
\[
w=\rho e^{i\theta},\qquad \rho>0,
\]
entonces las soluciones de $z^n=w$ son
\[
\boxed{z_k=\rho^{1/n}e^{i(\theta+2k\pi)/n}},
\qquad k=0,1,\ldots,n-1.
\]
En forma trigonométrica:
\[
z_k=\rho^{1/n}\left[
\cos\left(\frac{\theta+2k\pi}{n}\right)
+i\sin\left(\frac{\theta+2k\pi}{n}\right)\right].
\]
Todas tienen módulo $\rho^{1/n}$ y los argumentos consecutivos difieren en
$2\pi/n$.


\section{Ejemplos desarrollados paso a paso}

\begin{ejemplo}
Resolver
\[
z^3=-8+8\sqrt3\,i.
\]

\textbf{Paso 1. Escribir el número del lado derecho en forma polar.}

Sea
\[
w=-8+8\sqrt3\,i.
\]
Su módulo es
\[
|w|=\sqrt{(-8)^2+(8\sqrt3)^2}
=\sqrt{64+192}
=\sqrt{256}=16.
\]
Como $\Re(w)<0$ e $\Im(w)>0$, $w$ se encuentra en el segundo cuadrante.
Además,
\[
\tan\theta=\frac{8\sqrt3}{-8}=-\sqrt3.
\]
El ángulo de referencia es $\pi/3$; en el segundo cuadrante:
\[
\theta=\pi-\frac{\pi}{3}=\frac{2\pi}{3}.
\]
Por tanto,
\[
w=16e^{i2\pi/3}.
\]

\textbf{Paso 2. Aplicar la fórmula de las raíces cúbicas.}

Si $z^3=w$, entonces
\[
z_k=16^{1/3}
e^{i(2\pi/3+2k\pi)/3},
\qquad k=0,1,2.
\]
Como
\[
16^{1/3}=2\sqrt[3]{2},
\]
obtenemos
\[
z_k=2\sqrt[3]{2}\,
e^{i(2\pi/9+2k\pi/3)}.
\]

\textbf{Paso 3. Enumerar las tres raíces.}
\[
\boxed{
\begin{aligned}
z_0&=2\sqrt[3]{2}\,e^{i2\pi/9},\\
z_1&=2\sqrt[3]{2}\,e^{i8\pi/9},\\
z_2&=2\sqrt[3]{2}\,e^{i14\pi/9}.
\end{aligned}}
\]
Las tres raíces tienen módulo $2\sqrt[3]{2}$ y están separadas angularmente por
$2\pi/3$.
\end{ejemplo}

\begin{ejemplo}
Resolver
\[
z^4=-16+16i.
\]

\textbf{Paso 1. Forma polar.}
\[
|-16+16i|=\sqrt{256+256}=16\sqrt2.
\]
El número está en el segundo cuadrante y su ángulo de referencia es $\pi/4$,
por lo que
\[
\theta=\frac{3\pi}{4}.
\]
Así,
\[
-16+16i=16\sqrt2\,e^{i3\pi/4}.
\]

\textbf{Paso 2. Extraer las cuatro raíces.}
\[
z_k=(16\sqrt2)^{1/4}
e^{i(3\pi/4+2k\pi)/4},
\qquad k=0,1,2,3.
\]
Como
\[
(16\sqrt2)^{1/4}=2\,2^{1/8},
\]
resulta
\[
z_k=2\,2^{1/8}e^{i(3\pi/16+k\pi/2)}.
\]

\textbf{Paso 3. Enumerar.}
\[
\boxed{
\begin{aligned}
z_0&=2\,2^{1/8}e^{i3\pi/16},\\
z_1&=2\,2^{1/8}e^{i11\pi/16},\\
z_2&=2\,2^{1/8}e^{i19\pi/16},\\
z_3&=2\,2^{1/8}e^{i27\pi/16}.
\end{aligned}}
\]
\end{ejemplo}

\begin{ejemplo}[Reducción mediante sustitución]
Resolver
\[
z^6+7z^3-8=0.
\]

\textbf{Paso 1. Reconocer la estructura.}
El polinomio es cuadrático en $z^3$. Definimos
\[
u=z^3.
\]
Entonces
\[
u^2+7u-8=0.
\]

\textbf{Paso 2. Resolver la ecuación auxiliar.}
\[
u^2+7u-8=(u-1)(u+8)=0,
\]
por lo que
\[
u=1\qquad\text{o}\qquad u=-8.
\]
Regresando a $z$:
\[
z^3=1
\qquad\text{o}\qquad
z^3=-8.
\]

\textbf{Paso 3. Resolver $z^3=1$.}
\[
1=e^{i2k\pi},
\]
de modo que
\[
z_k=e^{i2k\pi/3},\qquad k=0,1,2.
\]
Así,
\[
1,\qquad
-\frac12+\frac{\sqrt3}{2}i,\qquad
-\frac12-\frac{\sqrt3}{2}i.
\]

\textbf{Paso 4. Resolver $z^3=-8$.}
\[
-8=8e^{i(\pi+2k\pi)}.
\]
Por tanto,
\[
z_k=2e^{i(\pi+2k\pi)/3},\qquad k=0,1,2,
\]
y obtenemos
\[
1+\sqrt3\,i,\qquad -2,\qquad 1-\sqrt3\,i.
\]

\textbf{Conclusión.}
\[
\boxed{
1,\;
-\frac12+\frac{\sqrt3}{2}i,\;
-\frac12-\frac{\sqrt3}{2}i,\;
1+\sqrt3 i,\;
-2,\;
1-\sqrt3 i.
}
\]
Son seis raíces distintas, como corresponde al grado $6$.
\end{ejemplo}

\begin{ejemplo}[Polinomio bicuadrado]
Resolver
\[
z^4-6z^2+25=0.
\]

\textbf{Paso 1. Sustitución.}
Sea $u=z^2$. Entonces
\[
u^2-6u+25=0.
\]

\textbf{Paso 2. Fórmula general.}
\[
u=\frac{6\pm\sqrt{36-100}}{2}
=\frac{6\pm\sqrt{-64}}{2}
=3\pm4i.
\]

\textbf{Paso 3. Resolver $z^2=3+4i$.}
Buscamos $z=a+bi$. Entonces
\[
(a+bi)^2=(a^2-b^2)+2ab\,i=3+4i.
\]
Por comparación,
\[
a^2-b^2=3,\qquad 2ab=4.
\]
Una solución evidente es $a=2$, $b=1$:
\[
(2+i)^2=3+4i.
\]
Las dos raíces son
\[
2+i,\qquad -2-i.
\]

\textbf{Paso 4. Resolver $z^2=3-4i$.}
Análogamente,
\[
(2-i)^2=3-4i,
\]
por lo que las raíces son
\[
2-i,\qquad -2+i.
\]

\textbf{Conclusión.}
\[
\boxed{2+i,\;-2-i,\;2-i,\;-2+i.}
\]
Y la factorización completa es
\[
\boxed{
z^4-6z^2+25=
(z-(2+i))(z+(2+i))(z-(2-i))(z+(2-i)).
}
\]
\end{ejemplo}

\begin{ejemplo}[Raíces desplazadas]
Resolver
\[
(z-1)^4=-16.
\]

\textbf{Paso 1. Trasladar la variable.}
Sea
\[
w=z-1.
\]
Entonces
\[
w^4=-16.
\]

\textbf{Paso 2. Forma polar.}
\[
-16=16e^{i(\pi+2k\pi)}.
\]
Por tanto,
\[
w_k=2e^{i(\pi+2k\pi)/4}
=2e^{i(\pi/4+k\pi/2)}.
\]

\textbf{Paso 3. Obtener las cuatro raíces de $w$.}
\[
\begin{aligned}
w_0&=2e^{i\pi/4}=\sqrt2+i\sqrt2,\\
w_1&=2e^{i3\pi/4}=-\sqrt2+i\sqrt2,\\
w_2&=2e^{i5\pi/4}=-\sqrt2-i\sqrt2,\\
w_3&=2e^{i7\pi/4}=\sqrt2-i\sqrt2.
\end{aligned}
\]

\textbf{Paso 4. Regresar a $z$.}
Como $z=w+1$:
\[
\boxed{
\begin{aligned}
z_0&=1+\sqrt2+i\sqrt2,\\
z_1&=1-\sqrt2+i\sqrt2,\\
z_2&=1-\sqrt2-i\sqrt2,\\
z_3&=1+\sqrt2-i\sqrt2.
\end{aligned}}
\]
Geométricamente, las cuatro raíces están sobre la circunferencia
\[
|z-1|=2.
\]
\end{ejemplo}

\section{Batería de ejercicios resueltos para clase}

\subsection{Nivel I: forma polar y raíces $n$-ésimas}

\begin{ejercicio}
Resolver
\[
z^3=-8+8i.
\]
\end{ejercicio}

\textbf{Solución.}
El módulo es
\[
\rho=\sqrt{(-8)^2+8^2}=8\sqrt2,
\]
y un argumento es
\[
\theta=\frac{3\pi}{4}.
\]
Así,
\[
-8+8i=8\sqrt2\,e^{i3\pi/4}.
\]
Por tanto,
\[
z_k=(8\sqrt2)^{1/3}
e^{i(3\pi/4+2k\pi)/3},\qquad k=0,1,2.
\]
Como $(8\sqrt2)^{1/3}=2\,2^{1/6}$,
\[
\boxed{
z_k=2\,2^{1/6}e^{i(\pi/4+2k\pi/3)},\quad k=0,1,2.
}
\]

\begin{ejercicio}
Resolver
\[
z^4=16i.
\]
\end{ejercicio}

\textbf{Solución.}
\[
16i=16e^{i(\pi/2+2k\pi)}.
\]
Entonces
\[
z_k=2e^{i(\pi/8+k\pi/2)},\qquad k=0,1,2,3.
\]
Por tanto,
\[
\boxed{
2e^{i\pi/8},\;
2e^{i5\pi/8},\;
2e^{i9\pi/8},\;
2e^{i13\pi/8}.
}
\]

\begin{ejercicio}
Resolver
\[
z^5=-32i.
\]
\end{ejercicio}

\textbf{Solución.}
Podemos tomar
\[
-32i=32e^{-i\pi/2}.
\]
Así,
\[
z_k=2e^{i(-\pi/2+2k\pi)/5}
=2e^{i(-\pi/10+2k\pi/5)},
\quad k=0,\ldots,4.
\]
Por tanto,
\[
\boxed{
z_k=2e^{i(-\pi/10+2k\pi/5)},\quad k=0,\ldots,4.
}
\]

\begin{ejercicio}
Resolver
\[
z^4=8(1-i).
\]
\end{ejercicio}

\textbf{Solución.}
Primero,
\[
|8(1-i)|=8\sqrt2,\qquad \arg(1-i)=-\frac{\pi}{4}.
\]
Entonces
\[
8(1-i)=8\sqrt2\,e^{-i\pi/4}.
\]
Las raíces son
\[
\boxed{
z_k=(8\sqrt2)^{1/4}
e^{i(-\pi/16+k\pi/2)},\quad k=0,1,2,3.
}
\]

\begin{ejercicio}
Resolver
\[
z^6=-64.
\]
\end{ejercicio}

\textbf{Solución.}
\[
-64=64e^{i(\pi+2k\pi)}.
\]
Como $64^{1/6}=2$,
\[
z_k=2e^{i(\pi+2k\pi)/6}
=2e^{i(\pi/6+k\pi/3)},
\quad k=0,\ldots,5.
\]
Así,
\[
\boxed{
\sqrt3+i,\;
2i,\;
-\sqrt3+i,\;
-\sqrt3-i,\;
-2i,\;
\sqrt3-i.
}
\]

\begin{ejercicio}
Resolver
\[
z^5=16(1+\sqrt3\,i).
\]
\end{ejercicio}

\textbf{Solución.}
\[
|16(1+\sqrt3 i)|=16\sqrt{1+3}=32,
\]
y
\[
\arg(1+\sqrt3 i)=\frac{\pi}{3}.
\]
Por tanto,
\[
16(1+\sqrt3 i)=32e^{i\pi/3}.
\]
Como $32^{1/5}=2$,
\[
\boxed{
z_k=2e^{i(\pi/15+2k\pi/5)},\quad k=0,\ldots,4.
}
\]

\subsection{Nivel II: sustitución y estructura algebraica}

\begin{ejercicio}
Resolver
\[
z^4+8z^2+16=0.
\]
\end{ejercicio}

\textbf{Solución.}
\[
z^4+8z^2+16=(z^2+4)^2.
\]
Por tanto,
\[
z^2=-4,
\]
de donde
\[
z=\pm2i.
\]
Como el factor está al cuadrado,
\[
\boxed{2i\text{ y }-2i,\quad\text{cada una de multiplicidad }2.}
\]

\begin{ejercicio}
Resolver
\[
z^6-7z^3-8=0.
\]
\end{ejercicio}

\textbf{Solución.}
Sea $u=z^3$:
\[
u^2-7u-8=(u-8)(u+1)=0.
\]
Así,
\[
z^3=8\qquad\text{o}\qquad z^3=-1.
\]
Para $z^3=8$:
\[
z=2e^{2k\pi i/3},\quad k=0,1,2.
\]
Para $z^3=-1$:
\[
z=e^{i(\pi+2k\pi)/3},\quad k=0,1,2.
\]
Por tanto,
\[
\boxed{
2,\;-1+i\sqrt3,\;-1-i\sqrt3,\;
\frac12+\frac{\sqrt3}{2}i,\;-1,\;
\frac12-\frac{\sqrt3}{2}i.
}
\]

\begin{ejercicio}
Resolver
\[
z^8-17z^4+16=0.
\]
\end{ejercicio}

\textbf{Solución.}
Sea $u=z^4$:
\[
u^2-17u+16=(u-1)(u-16)=0.
\]
Luego
\[
z^4=1\qquad\text{o}\qquad z^4=16.
\]
Del primer caso:
\[
1,\ i,\ -1,\ -i.
\]
Del segundo:
\[
2,\ 2i,\ -2,\ -2i.
\]
Por tanto hay ocho raíces distintas:
\[
\boxed{\pm1,\;\pm i,\;\pm2,\;\pm2i.}
\]

\begin{ejercicio}
Resolver
\[
z^6+2z^3+5=0.
\]
\end{ejercicio}

\textbf{Solución.}
Sea $u=z^3$:
\[
u^2+2u+5=0.
\]
Entonces
\[
u=-1\pm2i.
\]
Para cada uno de estos dos números debemos calcular tres raíces cúbicas.

Sea
\[
-1+2i=\sqrt5\,e^{i\theta},
\qquad
\theta=\pi-\arctan 2.
\]
Entonces
\[
z_k=5^{1/6}e^{i(\theta+2k\pi)/3},
\quad k=0,1,2.
\]
Para el conjugado,
\[
-1-2i=\sqrt5\,e^{-i\theta},
\]
y
\[
z_k=5^{1/6}e^{i(-\theta+2k\pi)/3},
\quad k=0,1,2.
\]
Así se obtienen las seis raíces:
\[
\boxed{
5^{1/6}e^{i(\pm\theta+2k\pi)/3},
\quad k=0,1,2,\qquad
\theta=\pi-\arctan2.
}
\]

\begin{ejercicio}
Resolver
\[
z^8+2z^4+4=0.
\]
\end{ejercicio}

\textbf{Solución.}
Sea $u=z^4$:
\[
u^2+2u+4=0.
\]
Entonces
\[
u=-1\pm i\sqrt3.
\]
Estos números tienen módulo $2$ y argumentos $\pm2\pi/3$:
\[
u=2e^{\pm i2\pi/3}.
\]
Por tanto, para cada signo,
\[
z_k=2^{1/4}
e^{i(\pm2\pi/3+2k\pi)/4},
\qquad k=0,1,2,3.
\]
En total:
\[
\boxed{
z_k=2^{1/4}e^{i(\pm\pi/6+k\pi/2)},
\quad k=0,1,2,3.
}
\]

\begin{ejercicio}
Resolver
\[
z^{10}-31z^5-32=0.
\]
\end{ejercicio}

\textbf{Solución.}
Sea $u=z^5$:
\[
u^2-31u-32=(u-32)(u+1)=0.
\]
Así,
\[
z^5=32\qquad\text{o}\qquad z^5=-1.
\]
Del primer caso:
\[
z_k=2e^{2k\pi i/5},\quad k=0,\ldots,4.
\]
Del segundo:
\[
z_k=e^{i(\pi+2k\pi)/5},\quad k=0,\ldots,4.
\]
Por tanto:
\[
\boxed{
\{2e^{2k\pi i/5}\}_{k=0}^4
\cup
\{e^{i(\pi+2k\pi)/5}\}_{k=0}^4.
}
\]

\subsection{Nivel III: desplazamientos y composiciones algebraicas}

\begin{ejercicio}
Resolver
\[
(z-2)^3=8i.
\]
\end{ejercicio}

\textbf{Solución.}
Sea $w=z-2$. Como
\[
8i=8e^{i\pi/2},
\]
\[
w_k=2e^{i(\pi/6+2k\pi/3)},\quad k=0,1,2.
\]
Finalmente,
\[
\boxed{
z_k=2+2e^{i(\pi/6+2k\pi/3)},\quad k=0,1,2.
}
\]

\begin{ejercicio}
Resolver
\[
(z+i)^4=16.
\]
\end{ejercicio}

\textbf{Solución.}
Sea $w=z+i$. Las raíces cuartas de $16$ son
\[
2,\;2i,\;-2,\;-2i.
\]
Como $z=w-i$:
\[
\boxed{
2-i,\quad i,\quad -2-i,\quad -3i.
}
\]

\begin{ejercicio}
Resolver
\[
(z-1+i)^5=-32.
\]
\end{ejercicio}

\textbf{Solución.}
Sea
\[
w=z-1+i.
\]
Entonces
\[
w^5=-32=32e^{i(\pi+2k\pi)}.
\]
Así,
\[
w_k=2e^{i(\pi+2k\pi)/5}.
\]
Como
\[
z=w+1-i,
\]
resulta
\[
\boxed{
z_k=1-i+2e^{i(\pi+2k\pi)/5},
\quad k=0,\ldots,4.
}
\]

\begin{ejercicio}
Resolver
\[
(z+2)^6+64=0.
\]
\end{ejercicio}

\textbf{Solución.}
\[
(z+2)^6=-64.
\]
Sea $w=z+2$. Entonces
\[
w_k=2e^{i(\pi/6+k\pi/3)},\quad k=0,\ldots,5.
\]
Por tanto,
\[
\boxed{
z_k=-2+2e^{i(\pi/6+k\pi/3)},\quad k=0,\ldots,5.
}
\]

\begin{ejercicio}
Resolver
\[
\bigl((z-1)^2\bigr)^2+4=0.
\]
\end{ejercicio}

\textbf{Solución.}
\[
(z-1)^4=-4.
\]
Como
\[
-4=4e^{i(\pi+2k\pi)},
\]
\[
z-1=\sqrt2\,e^{i(\pi/4+k\pi/2)},\quad k=0,1,2,3.
\]
Así,
\[
\boxed{
z_k=1+\sqrt2\,e^{i(\pi/4+k\pi/2)},\quad k=0,1,2,3.
}
\]

\begin{ejercicio}
Resolver
\[
(z^2+1)^2+16=0.
\]
\end{ejercicio}

\textbf{Solución.}
\[
(z^2+1)^2=-16.
\]
Por tanto,
\[
z^2+1=\pm4i,
\]
y entonces
\[
z^2=-1+4i
\qquad\text{o}\qquad
z^2=-1-4i.
\]
Para $z^2=a+bi$, una forma explícita de la raíz cuadrada es
\[
\sqrt{a+bi}
=
\sqrt{\frac{r+a}{2}}
+i\,\operatorname{sgn}(b)
\sqrt{\frac{r-a}{2}},
\quad r=\sqrt{a^2+b^2}.
\]
Aquí $r=\sqrt{17}$. Por tanto las cuatro raíces son
\[
\boxed{
\pm\left(
\sqrt{\frac{\sqrt{17}-1}{2}}
+i\sqrt{\frac{\sqrt{17}+1}{2}}
\right)
}
\]
y
\[
\boxed{
\pm\left(
\sqrt{\frac{\sqrt{17}-1}{2}}
-i\sqrt{\frac{\sqrt{17}+1}{2}}
\right).
}
\]

\section{Teorema Fundamental del Álgebra}

\begin{teorema}[Teorema Fundamental del Álgebra]
Sea
\[
p(z)=a_nz^n+a_{n-1}z^{n-1}+\cdots+a_1z+a_0,
\qquad a_n\neq0,
\]
un polinomio no constante con coeficientes complejos. Entonces existe
$z_0\in\C$ tal que
\[
p(z_0)=0.
\]
\end{teorema}

\begin{proposicion}
Todo polinomio de grado $n\ge1$ con coeficientes complejos puede factorizarse
como
\[
\boxed{p(z)=a_n(z-z_1)\cdots(z-z_n),}
\]
donde las raíces se cuentan de acuerdo con sus multiplicidades.
\end{proposicion}

\begin{proof}
Por el TFA, $p$ posee al menos una raíz $z_1$. El teorema del factor implica
\[
p(z)=(z-z_1)q_{n-1}(z).
\]
Si $n-1\ge1$, se aplica nuevamente el TFA a $q_{n-1}$. Repitiendo el
procedimiento hasta llegar a un polinomio constante se obtiene
\[
p(z)=a_n(z-z_1)\cdots(z-z_n).
\]
\end{proof}

\begin{observacion}
Un polinomio de grado $n$ tiene exactamente $n$ raíces en $\C$ cuando se
cuentan multiplicidades. Puede tener menos de $n$ raíces \emph{distintas}.
\end{observacion}

\begin{proposicion}[Raíces conjugadas]
Si $p$ tiene coeficientes reales y $z_0$ es raíz de $p$, entonces
$\overline{z_0}$ también es raíz, con la misma multiplicidad.
\end{proposicion}

\begin{proof}
Si los coeficientes son reales,
\[
p(\overline z)=\overline{p(z)}.
\]
Si $p(z_0)=0$,
\[
p(\overline{z_0})
=\overline{p(z_0)}
=0.
\]
La afirmación sobre multiplicidad se obtiene conjugando la factorización
correspondiente.
\end{proof}

\section{Ejemplos del TFA desarrollados}

\begin{ejemplo}[Cuatro raíces distintas]
Sea
\[
p(z)=z^4-16.
\]
Su grado es $4$. Factorizamos:
\[
z^4-16=(z^2-4)(z^2+4)
=(z-2)(z+2)(z-2i)(z+2i).
\]
Por tanto,
\[
\boxed{2,\;-2,\;2i,\;-2i}
\]
son cuatro raíces simples. La suma de multiplicidades es
\[
1+1+1+1=4.
\]
\end{ejemplo}

\begin{ejemplo}[Raíces múltiples]
\[
p(z)=z^4-2z^2+1=(z^2-1)^2=(z-1)^2(z+1)^2.
\]
Las raíces distintas son $1$ y $-1$, pero
\[
m(1)=2,\qquad m(-1)=2.
\]
Así,
\[
2+2=4=\deg p.
\]
\end{ejemplo}

\begin{ejemplo}[Raíces no reales]
\[
p(z)=z^4+5z^2+4.
\]
Sea $u=z^2$:
\[
u^2+5u+4=(u+1)(u+4)=0.
\]
Entonces
\[
z^2=-1\quad\text{o}\quad z^2=-4,
\]
por lo que
\[
\boxed{i,\;-i,\;2i,\;-2i.}
\]
Ninguna raíz es real, aunque el polinomio tiene coeficientes reales.
\end{ejemplo}

\begin{ejemplo}[Aplicación del teorema de raíces conjugadas]
Sea
\[
p(z)=z^3-3z^2+4z-2.
\]
Probamos raíces reales sencillas:
\[
p(1)=1-3+4-2=0.
\]
Luego $(z-1)$ es factor. Dividiendo:
\[
p(z)=(z-1)(z^2-2z+2).
\]
Resolvemos
\[
z^2-2z+2=0:
\]
\[
z=\frac{2\pm\sqrt{4-8}}2
=1\pm i.
\]
Por tanto,
\[
\boxed{
p(z)=(z-1)(z-(1+i))(z-(1-i)).
}
\]
La pareja $1\pm i$ ilustra el teorema de las raíces conjugadas.
\end{ejemplo}

\section{Ejercicios resueltos sobre el TFA}

En cada caso determinaremos grado, raíces, multiplicidades y factorización.

\begin{ejercicio}
\[
p(z)=z^3-8.
\]
\end{ejercicio}

\textbf{Solución.}
El grado es $3$. Resolver $z^3=8$ da
\[
z_k=2e^{2k\pi i/3},\quad k=0,1,2.
\]
Así,
\[
2,\quad -1+i\sqrt3,\quad -1-i\sqrt3.
\]
Todas son simples y
\[
\boxed{
p(z)=(z-2)(z+1-i\sqrt3)(z+1+i\sqrt3).
}
\]

\begin{ejercicio}
\[
p(z)=z^4+4.
\]
\end{ejercicio}

\textbf{Solución.}
\[
z^4=-4=4e^{i(\pi+2k\pi)}.
\]
Por tanto,
\[
z_k=\sqrt2\,e^{i(\pi/4+k\pi/2)},\quad k=0,1,2,3.
\]
En forma binómica:
\[
1+i,\quad -1+i,\quad -1-i,\quad 1-i.
\]
Todas son simples:
\[
\boxed{
p(z)=(z-1-i)(z+1-i)(z+1+i)(z-1+i).
}
\]

\begin{ejercicio}
\[
p(z)=z^4+2z^2+1.
\]
\end{ejercicio}

\textbf{Solución.}
\[
p(z)=(z^2+1)^2=(z-i)^2(z+i)^2.
\]
Las raíces son
\[
\boxed{i\text{ y }-i}
\]
con multiplicidad $2$ cada una. La suma de multiplicidades es $4$.

\begin{ejercicio}
\[
p(z)=z^4+13z^2+36.
\]
\end{ejercicio}

\textbf{Solución.}
Sea $u=z^2$:
\[
u^2+13u+36=(u+4)(u+9)=0.
\]
Entonces
\[
z^2=-4\quad\text{o}\quad z^2=-9.
\]
Las raíces son
\[
\boxed{2i,-2i,3i,-3i}.
\]
Todas son simples y
\[
\boxed{
p(z)=(z-2i)(z+2i)(z-3i)(z+3i).
}
\]

\begin{ejercicio}
\[
p(z)=z^6-1.
\]
\end{ejercicio}

\textbf{Solución.}
Las raíces sextas de la unidad son
\[
z_k=e^{ik\pi/3},\quad k=0,\ldots,5.
\]
Es decir,
\[
1,\quad
\frac12+\frac{\sqrt3}{2}i,\quad
-\frac12+\frac{\sqrt3}{2}i,\quad
-1,\quad
-\frac12-\frac{\sqrt3}{2}i,\quad
\frac12-\frac{\sqrt3}{2}i.
\]
Son seis raíces simples, de modo que
\[
\boxed{
z^6-1=\prod_{k=0}^{5}(z-e^{ik\pi/3}).
}
\]

\begin{ejercicio}
\[
p(z)=z^6+7z^3-8.
\]
\end{ejercicio}

\textbf{Solución.}
Este ejercicio coincide con uno de los ejemplos iniciales. Con $u=z^3$:
\[
(u-1)(u+8)=0.
\]
Las seis raíces son
\[
1,\;
-\frac12\pm\frac{\sqrt3}{2}i,\;
-2,\;
1\pm\sqrt3 i.
\]
Todas son simples; por tanto la suma de multiplicidades es $6$.

\begin{ejercicio}
\[
p(z)=z^8-17z^4+16.
\]
\end{ejercicio}

\textbf{Solución.}
Con $u=z^4$:
\[
(u-1)(u-16)=0.
\]
Así,
\[
z^4=1\quad\text{o}\quad z^4=16.
\]
Las raíces son
\[
\boxed{\pm1,\;\pm i,\;\pm2,\;\pm2i.}
\]
Son ocho raíces simples, como exige el grado $8$.

\begin{ejercicio}
\[
p(z)=(z-2)^3(z^2+4)^2.
\]
\end{ejercicio}

\textbf{Solución.}
El grado es
\[
3+2(2)=7.
\]
La raíz $z=2$ tiene multiplicidad $3$. Además,
\[
z^2+4=(z-2i)(z+2i),
\]
y al estar al cuadrado:
\[
z=2i,\quad z=-2i
\]
tienen multiplicidad $2$ cada una. La suma es
\[
3+2+2=7.
\]
La factorización completa ya es
\[
\boxed{
p(z)=(z-2)^3(z-2i)^2(z+2i)^2.
}
\]

\begin{ejercicio}
\[
p(z)=z^5-5z^4+10z^3-10z^2+5z-1.
\]
\end{ejercicio}

\textbf{Solución.}
Reconocemos el desarrollo binomial:
\[
(z-1)^5
=z^5-5z^4+10z^3-10z^2+5z-1.
\]
Por tanto,
\[
\boxed{z=1}
\]
es la única raíz distinta y tiene multiplicidad $5$:
\[
\boxed{p(z)=(z-1)^5.}
\]

\begin{ejercicio}
\[
p(z)=z^6+3z^4+3z^2+1.
\]
\end{ejercicio}

\textbf{Solución.}
Reconocemos
\[
z^6+3z^4+3z^2+1=(z^2+1)^3.
\]
Luego
\[
p(z)=(z-i)^3(z+i)^3.
\]
Por tanto,
\[
\boxed{i\text{ y }-i}
\]
tienen multiplicidad $3$ cada una, y
\[
3+3=6.
\]

\section{Ejercicios de razonamiento resueltos}

\begin{ejercicio}
¿Puede existir un polinomio de grado $5$ con coeficientes complejos que no
tenga ninguna raíz compleja?
\end{ejercicio}

\textbf{Respuesta.}
No. El TFA afirma que todo polinomio no constante con coeficientes complejos
tiene al menos una raíz en $\C$. De hecho, contando multiplicidades, un
polinomio de grado $5$ tiene exactamente cinco raíces complejas.

\begin{ejercicio}
Un polinomio de grado $7$ tiene exactamente tres raíces distintas. ¿Contradice
esto el TFA?
\end{ejercicio}

\textbf{Respuesta.}
No. Las tres raíces distintas pueden tener multiplicidades cuya suma sea $7$.
Por ejemplo,
\[
p(z)=(z-a)^3(z-b)^2(z-c)^2
\]
tiene grado $7$, sólo tres raíces distintas y multiplicidades
\[
3+2+2=7.
\]

\begin{ejercicio}
¿Puede un polinomio de grado impar con coeficientes reales tener únicamente
raíces no reales?
\end{ejercicio}

\textbf{Respuesta.}
No. Las raíces no reales de un polinomio real aparecen en pares conjugados.
Cada par aporta un número par de raíces, contando multiplicidades. Un grado
impar no puede componerse exclusivamente de pares; por ello debe existir al
menos una raíz real.

\begin{ejercicio}
Si $p$ tiene coeficientes reales y $2+3i$ es raíz de multiplicidad $2$, ¿qué
puede afirmarse acerca de $2-3i$?
\end{ejercicio}

\textbf{Respuesta.}
Por el teorema de raíces conjugadas,
\[
\boxed{2-3i}
\]
también es raíz y tiene la misma multiplicidad $2$.

\begin{ejercicio}
Construir un polinomio mónico real de grado $4$ cuyas raíces sean
\[
1+i,\quad1-i,\quad-2+3i,\quad-2-3i.
\]
\end{ejercicio}

\textbf{Solución.}
Agrupamos conjugados:
\[
(z-(1+i))(z-(1-i))
=((z-1)-i)((z-1)+i)
=(z-1)^2+1
=z^2-2z+2.
\]
Asimismo,
\[
(z+2-3i)(z+2+3i)
=(z+2)^2+9
=z^2+4z+13.
\]
Entonces
\[
p(z)=(z^2-2z+2)(z^2+4z+13).
\]
Multiplicando:
\[
\boxed{
p(z)=z^4+2z^3+7z^2-18z+26.
}
\]

\begin{ejercicio}
Construir un polinomio mónico real de grado $5$ que tenga $2+i$ como raíz
doble y $-1$ como raíz simple.
\end{ejercicio}

\textbf{Solución.}
Como los coeficientes deben ser reales, $2-i$ también debe ser raíz doble.
Ya tenemos
\[
2+2+1=5
\]
raíces contando multiplicidades, exactamente el grado requerido. Por tanto,
\[
p(z)=(z-(2+i))^2(z-(2-i))^2(z+1).
\]
El producto del par conjugado es
\[
(z-(2+i))(z-(2-i))
=(z-2)^2+1
=z^2-4z+5.
\]
Luego
\[
p(z)=(z^2-4z+5)^2(z+1).
\]
Como
\[
(z^2-4z+5)^2=z^4-8z^3+26z^2-40z+25,
\]
obtenemos
\[
\boxed{
p(z)=z^5-7z^4+18z^3-14z^2-15z+25.
}
\]

\begin{ejercicio}
¿Por qué $z^2+1$ muestra que $\mathbb{R}$ no es algebraicamente cerrado?
\end{ejercicio}

\textbf{Respuesta.}
El polinomio
\[
p(x)=x^2+1
\]
tiene coeficientes reales, pero la ecuación
\[
x^2+1=0
\]
implica
\[
x^2=-1,
\]
que no tiene solución en $\mathbb{R}$. Por ello existe un polinomio no
constante con coeficientes reales que no tiene raíces reales. En consecuencia,
$\mathbb{R}$ no es algebraicamente cerrado.

\begin{ejercicio}
¿En qué sentido el TFA afirma que $\C$ es algebraicamente cerrado?
\end{ejercicio}

\textbf{Respuesta.}
Un campo es algebraicamente cerrado si todo polinomio no constante con
coeficientes en ese campo posee una raíz en el mismo campo. El TFA afirma
precisamente que todo polinomio no constante con coeficientes complejos posee
una raíz compleja. Aplicando repetidamente el teorema del factor, todo
polinomio se descompone completamente en factores lineales sobre $\C$.

\section{Problemas adicionales resueltos}

\begin{ejercicio}
Resolver
\[
z^4=81\left(\cos\frac{2\pi}{3}+i\sin\frac{2\pi}{3}\right).
\]
\end{ejercicio}

\textbf{Solución.}
El módulo de cada raíz es
\[
81^{1/4}=3.
\]
Los argumentos son
\[
\frac{2\pi/3+2k\pi}{4}
=\frac{\pi}{6}+\frac{k\pi}{2}.
\]
Por tanto,
\[
\boxed{
z_k=3e^{i(\pi/6+k\pi/2)},\quad k=0,1,2,3.
}
\]

\begin{ejercicio}
Resolver
\[
z^5=32\left(\cos\frac{5\pi}{6}+i\sin\frac{5\pi}{6}\right).
\]
\end{ejercicio}

\textbf{Solución.}
Como $32^{1/5}=2$,
\[
\arg z_k=
\frac{5\pi/6+2k\pi}{5}
=\frac{\pi}{6}+\frac{2k\pi}{5}.
\]
Así,
\[
\boxed{
z_k=2e^{i(\pi/6+2k\pi/5)},\quad k=0,\ldots,4.
}
\]

\begin{ejercicio}
Determinar todas las raíces de
\[
z^8+15z^4-16=0.
\]
\end{ejercicio}

\textbf{Solución.}
Sea $u=z^4$:
\[
u^2+15u-16=(u-1)(u+16)=0.
\]
Entonces
\[
z^4=1
\quad\text{o}\quad
z^4=-16.
\]
Del primer caso:
\[
1,\;i,\;-1,\;-i.
\]
Del segundo:
\[
z_k=2e^{i(\pi/4+k\pi/2)},\quad k=0,1,2,3,
\]
es decir,
\[
\sqrt2(1+i),\quad
\sqrt2(-1+i),\quad
\sqrt2(-1-i),\quad
\sqrt2(1-i).
\]
En total hay ocho raíces distintas.

\begin{ejercicio}
Determinar todas las raíces de
\[
z^9-7z^6+14z^3-8=0.
\]
\end{ejercicio}

\textbf{Solución.}
Sea $u=z^3$. Entonces
\[
u^3-7u^2+14u-8=0.
\]
Probamos $u=1$:
\[
1-7+14-8=0.
\]
Por tanto $(u-1)$ es factor. Dividiendo:
\[
u^3-7u^2+14u-8=(u-1)(u^2-6u+8)
=(u-1)(u-2)(u-4).
\]
Así,
\[
z^3=1,\qquad z^3=2,\qquad z^3=4.
\]
Por tanto las nueve raíces son
\[
\boxed{
e^{2k\pi i/3},\quad
2^{1/3}e^{2k\pi i/3},\quad
4^{1/3}e^{2k\pi i/3},
\qquad k=0,1,2.
}
\]

\begin{ejercicio}
Factorizar completamente sobre $\C$:
\[
z^6+3z^4+4z^2+12.
\]
\end{ejercicio}

\textbf{Solución.}
Agrupamos:
\[
z^6+3z^4+4z^2+12
=z^4(z^2+3)+4(z^2+3)
=(z^2+3)(z^4+4).
\]
Ahora,
\[
z^2+3=0
\quad\Longrightarrow\quad
z=\pm i\sqrt3.
\]
Además,
\[
z^4=-4,
\]
cuyas raíces son
\[
1+i,\quad -1+i,\quad -1-i,\quad 1-i.
\]
Por tanto,
\[
\boxed{
\begin{aligned}
z^6+3z^4+4z^2+12
={}&(z-i\sqrt3)(z+i\sqrt3)\\
&\cdot(z-1-i)(z+1-i)(z+1+i)(z-1+i).
\end{aligned}}
\]

\begin{ejercicio}
Construir un polinomio real de grado $6$ que tenga como raíces
\[
1+i,\qquad -2+3i,\qquad4i,
\]
junto con todas las raíces que necesariamente deben acompañarlas.
\end{ejercicio}

\textbf{Solución.}
Por conjugación también deben ser raíces
\[
1-i,\qquad -2-3i,\qquad -4i.
\]
Así obtenemos exactamente seis raíces. El polinomio mónico es
\[
p(z)=
(z-(1+i))(z-(1-i))
(z+2-3i)(z+2+3i)
(z-4i)(z+4i).
\]
Agrupando conjugados:
\[
p(z)=(z^2-2z+2)(z^2+4z+13)(z^2+16).
\]
Ésta es ya una factorización real completa en factores cuadráticos
irreducibles. Si se desea expandir:
\[
(z^2-2z+2)(z^2+4z+13)
=z^4+2z^3+7z^2-18z+26,
\]
y finalmente
\[
\boxed{
p(z)=z^6+2z^5+23z^4+14z^3+138z^2-288z+416.
}
\]

\begin{ejercicio}
Un polinomio mónico real de grado $6$ tiene $1+2i$ como raíz doble y $-3$
como raíz doble. Construir el polinomio.
\end{ejercicio}

\textbf{Solución.}
Como $1+2i$ es raíz doble y los coeficientes son reales, $1-2i$ también es
raíz doble. Junto con $-3$ de multiplicidad $2$, tenemos
\[
2+2+2=6,
\]
por lo que ya se han determinado todas las raíces.

Así,
\[
p(z)=(z-(1+2i))^2(z-(1-2i))^2(z+3)^2.
\]
Agrupando conjugados:
\[
(z-(1+2i))(z-(1-2i))
=(z-1)^2+4
=z^2-2z+5.
\]
Por tanto,
\[
\boxed{
p(z)=(z^2-2z+5)^2(z+3)^2.
}
\]
Esta forma factorizada muestra inmediatamente las raíces y sus
multiplicidades.


\section{Logaritmo complejo}

\subsection{Motivación: por qué el logaritmo complejo es multivaluado}

En variable real, para $x>0$, el logaritmo natural $\ln x$ es la función
inversa de la exponencial real. En el plano complejo aparece una diferencia
fundamental: la exponencial compleja es periódica,
\[
e^{z+2\pi i k}=e^z,\qquad k\in\mathbb Z.
\]
Por esta razón, la ecuación
\[
e^w=z,\qquad z\neq0,
\]
no tiene una única solución $w$, sino infinitas.

Sea
\[
z=re^{i\theta},\qquad r=|z|>0.
\]
Buscamos $w=u+iv$ tal que
\[
e^w=z.
\]
Usando
\[
e^{u+iv}=e^u(\cos v+i\sin v),
\]
y comparando con
\[
z=r(\cos\theta+i\sin\theta),
\]
se obtienen las condiciones
\[
e^u=r
\qquad\text{y}\qquad
v=\theta+2k\pi,\qquad k\in\mathbb Z.
\]
De la primera,
\[
u=\ln r.
\]
Por tanto, todos los valores del logaritmo complejo son
\[
\boxed{
\log z=\ln|z|+i(\arg z+2k\pi),\qquad k\in\mathbb Z.
}
\]

Si $\theta$ representa un argumento particular de $z$, también escribimos
\[
\boxed{
\log z=\ln r+i(\theta+2k\pi),\qquad k\in\mathbb Z.
}
\]

\begin{observacion}
El símbolo $\log z$ se utiliza aquí para representar el logaritmo complejo
multivaluado. No debe confundirse con una rama univaluada del logaritmo.
\end{observacion}

\subsection{Argumento principal y logaritmo principal}

Como un número complejo no nulo tiene infinitos argumentos,
\[
\arg z=\theta+2k\pi,\qquad k\in\mathbb Z,
\]
es útil seleccionar uno de ellos.

Definimos el \emph{argumento principal} mediante
\[
\operatorname{Arg}z\in(-\pi,\pi].
\]
Entonces el \emph{logaritmo principal} se define por
\[
\boxed{
\operatorname{Log}z=\ln|z|+i\operatorname{Arg}z.
}
\]

Por tanto,
\[
\log z
=
\operatorname{Log}z+2k\pi i,
\qquad k\in\mathbb Z,
\]
siempre que se interprete adecuadamente el argumento principal.

\begin{observacion}
El número $z=0$ no posee logaritmo complejo. En efecto, no existe
$w\in\mathbb C$ tal que $e^w=0$, pues
\[
|e^w|=e^{\Re(w)}>0.
\]
\end{observacion}

\subsection{Ejemplos básicos completamente desarrollados}

\begin{ejemplo}
Calcular todos los valores de
\[
\log 1.
\]

\textbf{Paso 1. Forma polar.}

El módulo es
\[
|1|=1,
\]
y todos sus argumentos son
\[
2k\pi,\qquad k\in\mathbb Z.
\]

\textbf{Paso 2. Aplicar la definición.}
\[
\log1=\ln1+i(2k\pi).
\]
Como $\ln1=0$,
\[
\boxed{
\log1=2k\pi i,\qquad k\in\mathbb Z.
}
\]

El valor principal corresponde a $k=0$:
\[
\boxed{\operatorname{Log}1=0.}
\]

\textbf{Comprobación.}
\[
e^{2k\pi i}
=\cos(2k\pi)+i\sin(2k\pi)
=1.
\]
\end{ejemplo}

\begin{ejemplo}
Calcular todos los valores de
\[
\log(-1).
\]

\textbf{Paso 1. Módulo y argumento.}
\[
|-1|=1.
\]
Un argumento es $\pi$, por lo que todos los argumentos son
\[
\pi+2k\pi=(2k+1)\pi.
\]

\textbf{Paso 2. Sustituir.}
\[
\log(-1)
=\ln1+i(2k+1)\pi.
\]
Así,
\[
\boxed{
\log(-1)=(2k+1)\pi i,\qquad k\in\mathbb Z.
}
\]

Como $\operatorname{Arg}(-1)=\pi$,
\[
\boxed{\operatorname{Log}(-1)=i\pi.}
\]
\end{ejemplo}

\begin{ejemplo}
Calcular todos los valores de
\[
\log i.
\]

Tenemos
\[
|i|=1,\qquad
\arg i=\frac{\pi}{2}+2k\pi.
\]
Entonces
\[
\boxed{
\log i=i\left(\frac{\pi}{2}+2k\pi\right),
\qquad k\in\mathbb Z.
}
\]
El valor principal es
\[
\boxed{\operatorname{Log}i=\frac{\pi i}{2}.}
\]
\end{ejemplo}

\begin{ejemplo}
Calcular todos los valores de
\[
\log(1+i).
\]

\textbf{Paso 1. Módulo.}
\[
|1+i|=\sqrt{1^2+1^2}=\sqrt2.
\]

\textbf{Paso 2. Argumento.}

El número está en el primer cuadrante y
\[
\tan\theta=1,
\]
por lo que
\[
\theta=\frac{\pi}{4}.
\]

\textbf{Paso 3. Logaritmo.}
\[
\log(1+i)
=
\ln\sqrt2
+i\left(\frac{\pi}{4}+2k\pi\right).
\]
Como
\[
\ln\sqrt2=\frac12\ln2,
\]
obtenemos
\[
\boxed{
\log(1+i)
=
\frac12\ln2
+i\left(\frac{\pi}{4}+2k\pi\right),
\quad k\in\mathbb Z.
}
\]
El valor principal es
\[
\boxed{
\operatorname{Log}(1+i)
=
\frac12\ln2+\frac{\pi}{4}i.
}
\]
\end{ejemplo}

\begin{ejemplo}
Calcular todos los valores de
\[
\log(-1+i\sqrt3).
\]

\textbf{Paso 1. Módulo.}
\[
r=\sqrt{(-1)^2+(\sqrt3)^2}=2.
\]

\textbf{Paso 2. Argumento.}

El punto está en el segundo cuadrante. Su ángulo de referencia es $\pi/3$,
de modo que
\[
\theta=\frac{2\pi}{3}.
\]

\textbf{Paso 3. Logaritmo.}
\[
\boxed{
\log(-1+i\sqrt3)
=
\ln2+i\left(\frac{2\pi}{3}+2k\pi\right),
\quad k\in\mathbb Z.
}
\]
Por tanto,
\[
\boxed{
\operatorname{Log}(-1+i\sqrt3)
=
\ln2+\frac{2\pi}{3}i.
}
\]
\end{ejemplo}

\subsection{Ejemplos con determinación cuidadosa del argumento}

\begin{ejemplo}
Calcular $\log(-2-2i)$.

El módulo es
\[
|-2-2i|=\sqrt8=2\sqrt2.
\]
El punto está en el tercer cuadrante. Un argumento puede escribirse como
\[
\theta=\frac{5\pi}{4},
\]
pero para el argumento principal debemos escoger el representante en
$(-\pi,\pi]$:
\[
\operatorname{Arg}(-2-2i)=-\frac{3\pi}{4}.
\]
Por tanto,
\[
\boxed{
\log(-2-2i)
=
\ln(2\sqrt2)
+i\left(-\frac{3\pi}{4}+2k\pi\right),
\quad k\in\mathbb Z.
}
\]
Como
\[
\ln(2\sqrt2)=\frac32\ln2,
\]
el valor principal es
\[
\boxed{
\operatorname{Log}(-2-2i)
=
\frac32\ln2-\frac{3\pi}{4}i.
}
\]
\end{ejemplo}

\begin{ejemplo}
Calcular $\log(\sqrt3-i)$.

Tenemos
\[
|\sqrt3-i|=\sqrt{3+1}=2.
\]
El punto está en el cuarto cuadrante y
\[
\tan\theta=-\frac{1}{\sqrt3}.
\]
Por tanto,
\[
\operatorname{Arg}(\sqrt3-i)=-\frac{\pi}{6}.
\]
Así,
\[
\boxed{
\log(\sqrt3-i)
=
\ln2+i\left(-\frac{\pi}{6}+2k\pi\right),
\quad k\in\mathbb Z.
}
\]
y
\[
\boxed{
\operatorname{Log}(\sqrt3-i)=\ln2-\frac{\pi}{6}i.
}
\]
\end{ejemplo}

\subsection{Propiedades y precauciones}

Para $z_1,z_2\neq0$, los logaritmos multivaluados satisfacen, en el sentido
apropiado de conjuntos de valores,
\[
\log(z_1z_2)=\log z_1+\log z_2.
\]
Sin embargo, para el \emph{logaritmo principal} no siempre es cierto que
\[
\operatorname{Log}(z_1z_2)
=
\operatorname{Log}z_1+\operatorname{Log}z_2.
\]

\begin{ejemplo}[Una identidad que falla para el logaritmo principal]
Tomemos
\[
z_1=z_2=-1.
\]
Entonces
\[
z_1z_2=1
\]
y
\[
\operatorname{Log}(1)=0.
\]
Pero
\[
\operatorname{Log}(-1)=i\pi,
\]
por lo que
\[
\operatorname{Log}(-1)+\operatorname{Log}(-1)
=2\pi i.
\]
Así,
\[
\boxed{
\operatorname{Log}((-1)(-1))
\neq
\operatorname{Log}(-1)+\operatorname{Log}(-1).
}
\]
La diferencia es precisamente un múltiplo entero de $2\pi i$.
\end{ejemplo}

\begin{observacion}
Ésta es una de las razones por las que en variable compleja las identidades
algebraicas que involucran logaritmos deben manejarse con cuidado. La elección
de una rama afecta el valor obtenido.
\end{observacion}

\subsection{Ramas del logaritmo: el problema que debemos resolver}

La noción de \emph{rama} del logaritmo complejo merece una discusión cuidadosa.
No se trata simplemente de una convención de notación: aparece porque el
logaritmo complejo, considerado inicialmente, no es una función univaluada.

La idea central puede resumirse en
\[
\boxed{
\text{el problema no está en }\ln|z|;
\qquad
\text{el problema está en elegir el argumento de }z.
}
\]

En efecto, para $z\neq0$,
\[
\log z=\ln|z|+i(\theta+2k\pi),\qquad k\in\mathbb Z.
\]
El módulo $|z|$ es único y, por tanto, $\ln|z|$ también lo es. En cambio,
el argumento no es único:
\[
\theta,\quad\theta+2\pi,\quad\theta+4\pi,\quad\ldots
\]
representan la misma dirección en el plano complejo.

Así,
\[
\boxed{
\text{multivaluación de }\log z
\quad\Longleftrightarrow\quad
\text{multivaluación del argumento}.
}
\]

\subsubsection{¿Por qué necesitamos escoger una rama?}

Consideremos primero el número
\[
z=i.
\]
Podemos escribir
\[
i=e^{i\pi/2},
\]
pero también
\[
i=e^{i(\pi/2+2\pi)}
=e^{i(\pi/2-2\pi)}
=e^{i(\pi/2+2k\pi)}.
\]
Por consiguiente,
\[
\log i
=
i\left(\frac{\pi}{2}+2k\pi\right),
\qquad k\in\mathbb Z.
\]

Tenemos entonces infinitos valores asociados al mismo número $i$. Si queremos
que el logaritmo sea una \emph{función} en el sentido usual, debemos seleccionar
exactamente uno de ellos.

Una \emph{rama del argumento} consiste precisamente en escoger, de manera
coherente y continua sobre cierto dominio, un único argumento para cada número
complejo del dominio. Una vez elegida una rama del argumento, se obtiene una
rama del logaritmo mediante
\[
L(z)=\ln|z|+i\,\Theta(z),
\]
donde $\Theta(z)$ es la rama elegida del argumento.

\subsubsection{La elección principal}

Una elección particularmente importante consiste en tomar argumentos en un
intervalo de longitud $2\pi$ centrado alrededor de cero.

Como convención puntual, suele definirse el argumento principal mediante
\[
\operatorname{Arg}z\in(-\pi,\pi].
\]
Así, por ejemplo,
\[
\operatorname{Arg}(1)=0,\qquad
\operatorname{Arg}(i)=\frac{\pi}{2},\qquad
\operatorname{Arg}(-i)=-\frac{\pi}{2},\qquad
\operatorname{Arg}(-1)=\pi.
\]

Con esta convención podemos escribir el \emph{valor principal} del logaritmo:
\[
\boxed{
\operatorname{Log}z
=
\ln|z|+i\operatorname{Arg}z.
}
\]

Sin embargo, aquí es necesario hacer una distinción importante.

\begin{observacion}[Valor principal y rama principal analítica]
La convención
\[
\operatorname{Arg}z\in(-\pi,\pi]
\]
permite asignar un argumento principal incluso a los números reales negativos.
Pero si deseamos que el logaritmo sea una función continua y analítica en un
dominio abierto, debemos eliminar una semirrecta del plano.

Para la \emph{rama principal analítica} se trabaja usualmente en
\[
\boxed{
D=\mathbb C\setminus(-\infty,0]
}
\]
y se toma
\[
\boxed{
-\pi<\Arg z<\pi.
}
\]
Sobre este dominio se define
\[
\boxed{
\Log z=\ln|z|+i\Arg z.
}
\]
\end{observacion}

Esta diferencia entre ``valor principal'' y ``rama principal analítica'' es
importante: la primera es una selección de un valor; la segunda es una función
univaluada, continua y analítica sobre un dominio abierto apropiado.

\subsubsection{¿Es obligatoria la elección $(-\pi,\pi)$?}

No. La elección principal es una convención especialmente útil, pero no es la
única posible.

Podríamos, por ejemplo, escoger
\[
0<\theta<2\pi,
\]
lo cual obliga a colocar el corte sobre el eje real positivo. Más generalmente,
podemos seleccionar un intervalo
\[
\alpha<\theta<\alpha+2\pi.
\]
Esto conduce a una rama
\[
\Arg_\alpha z\in(\alpha,\alpha+2\pi)
\]
y a un logaritmo
\[
\Log_\alpha z
=
\ln|z|+i\Arg_\alpha z
\]
definido sobre el plano con la semirrecta de ángulo $\alpha$ eliminada.

Por tanto,
\[
\boxed{
\text{una rama es una elección coherente de una de las infinitas hojas
del argumento.}
}
\]

\subsubsection{¿Por qué es necesario hacer un corte?}

Éste es el punto geométrico fundamental.

Consideremos la circunferencia unitaria
\[
z=e^{i\theta}.
\]
Partimos de un punto cercano a $1$ y recorremos la circunferencia en sentido
antihorario. El argumento aumenta continuamente:
\[
0,\quad
\frac{\pi}{2},\quad
\pi,\quad
\frac{3\pi}{2},\quad
2\pi.
\]
Pero al llegar a $\theta=2\pi$ tenemos
\[
e^{i2\pi}=e^{i0}=1.
\]
Geométricamente hemos regresado exactamente al mismo punto, aunque el argumento
ha aumentado en $2\pi$.

Si continuamos otra vuelta,
\[
\theta=4\pi,
\]
y otra más,
\[
\theta=6\pi.
\]

Por consiguiente, si permitimos rodear completamente al origen, no podemos
mantener simultáneamente las tres propiedades siguientes:

\begin{enumerate}
\item asignar un único argumento a cada punto;
\item hacer que ese argumento varíe continuamente;
\item permitir recorridos cerrados que den una vuelta completa alrededor de
$0$.
\end{enumerate}

Para obtener una rama continua debemos impedir que el dominio permita completar
esa vuelta. Una forma de hacerlo es eliminar una semirrecta que parte del
origen.

Para la rama principal eliminamos
\[
(-\infty,0].
\]
Entonces el dominio es
\[
D=\mathbb C\setminus(-\infty,0].
\]

Ya no es posible recorrer una curva cerrada alrededor del origen sin atravesar
el corte.

\[
\boxed{
\text{El corte de rama impide completar una vuelta alrededor del origen
sin abandonar el dominio.}
}
\]

\subsubsection{Qué ocurre a ambos lados del corte}

La necesidad del corte puede verse de manera particularmente clara cerca del
eje real negativo.

Consideremos
\[
z_+=-1+i\varepsilon,
\qquad
z_-=-1-i\varepsilon,
\qquad \varepsilon>0.
\]

El punto $z_+$ se aproxima a $-1$ desde el semiplano superior. Por ello,
\[
\Arg(z_+)\longrightarrow\pi
\qquad
\text{cuando }\varepsilon\to0^+.
\]

En cambio, $z_-$ se aproxima a $-1$ desde el semiplano inferior:
\[
\Arg(z_-)\longrightarrow-\pi.
\]

Como
\[
|z_\pm|\longrightarrow1,
\]
tenemos
\[
\ln|z_\pm|\longrightarrow0.
\]
Así,
\[
\Log(z_+)\longrightarrow i\pi,
\]
mientras que
\[
\Log(z_-)\longrightarrow-i\pi.
\]

La diferencia entre los límites laterales es
\[
i\pi-(-i\pi)=2\pi i.
\]

Por tanto,
\[
\boxed{
\lim_{\substack{z\to-1\\\Im z>0}}\Log z=i\pi,
\qquad
\lim_{\substack{z\to-1\\\Im z<0}}\Log z=-i\pi.
}
\]

Ésta es una forma precisa de observar el salto de la rama principal a través
del corte:
\[
\boxed{\text{el salto es }2\pi i.}
\]

\begin{ejemplo}[Dos puntos muy próximos, dos argumentos muy diferentes]
Consideremos
\[
z_1=-1+0.01i,
\qquad
z_2=-1-0.01i.
\]
Ambos puntos están muy próximos en el plano complejo. Sin embargo,
\[
\Arg z_1\approx\pi,
\qquad
\Arg z_2\approx-\pi.
\]
En consecuencia,
\[
\Log z_1\approx i\pi,
\qquad
\Log z_2\approx-i\pi.
\]
Aunque $z_1$ y $z_2$ están arbitrariamente cerca cuando el término imaginario
tiende a cero, sus logaritmos principales se separan aproximadamente en
$2\pi i$. Esto muestra por qué no debemos intentar extender continuamente la
rama principal a través del eje real negativo.
\end{ejemplo}

\subsubsection{Interpretación geométrica mediante una vuelta alrededor del origen}

Podemos formalizar la intuición anterior tomando
\[
\gamma(t)=e^{it},
\qquad0\leq t\leq2\pi.
\]
Supongamos, por un momento, que existiera una función argumento continua
$\Theta$ definida a lo largo de toda la circunferencia y compatible con
\[
e^{i\Theta(\gamma(t))}=\gamma(t).
\]
Al recorrer la curva continuamente, necesariamente
\[
\Theta(\gamma(t))=\Theta(\gamma(0))+t
\]
salvo múltiplos constantes de $2\pi$. Al completar la vuelta:
\[
\Theta(\gamma(2\pi))
=
\Theta(\gamma(0))+2\pi.
\]
Pero
\[
\gamma(2\pi)=\gamma(0)=1.
\]
Una función univaluada tendría que satisfacer
\[
\Theta(\gamma(2\pi))
=
\Theta(\gamma(0)),
\]
lo cual contradice el incremento de $2\pi$.

Esta contradicción explica, de forma elemental, por qué no existe una rama
continua del argumento sobre todo
\[
\mathbb C\setminus\{0\}.
\]

\subsubsection{La rama principal como función analítica}

Una vez realizado el corte,
\[
D=\mathbb C\setminus(-\infty,0],
\]
podemos definir
\[
\Log z=\ln|z|+i\Arg z,
\qquad -\pi<\Arg z<\pi.
\]

Si
\[
z=x+iy,
\]
entonces
\[
\Log z
=
\frac12\ln(x^2+y^2)
+i\Arg(x+iy).
\]

En este dominio, $\Log z$ es una función analítica y satisface
\[
\boxed{
\frac{d}{dz}\Log z=\frac1z.
}
\]

Una manera de comprender esta identidad es recordar que, localmente,
\[
e^{\Log z}=z.
\]
Derivando formalmente:
\[
e^{\Log z}\frac{d}{dz}\Log z=1.
\]
Como
\[
e^{\Log z}=z,
\]
resulta
\[
z\frac{d}{dz}\Log z=1,
\]
y por tanto
\[
\boxed{
(\Log z)'=\frac1z.
}
\]

\subsubsection{Una consecuencia algebraica importante}

La elección de una rama también explica por qué ciertas identidades familiares
del logaritmo real pueden fallar para valores principales.

Por ejemplo,
\[
(-1)(-1)=1.
\]
Entonces
\[
\operatorname{Log}(1)=0.
\]
Sin embargo,
\[
\operatorname{Log}(-1)=i\pi,
\]
de modo que
\[
\operatorname{Log}(-1)+\operatorname{Log}(-1)=2\pi i.
\]
Así,
\[
\boxed{
\operatorname{Log}(z_1z_2)
\neq
\operatorname{Log}z_1+\operatorname{Log}z_2
}
\]
en general para valores principales.

No se trata de un fallo del álgebra, sino de que al tomar valores principales
estamos efectuando una selección particular entre valores que difieren por
múltiplos de $2\pi i$.

\subsubsection{Ejercicios resueltos sobre rama principal}

\begin{ejercicio}
Determine el valor principal de
\[
\log(-1-i).
\]
\end{ejercicio}

\textbf{Solución.}
El módulo es
\[
|-1-i|=\sqrt2.
\]
El punto está en el tercer cuadrante. Un argumento positivo es $5\pi/4$, pero
éste no pertenece a $(-\pi,\pi]$. Restamos $2\pi$:
\[
\frac{5\pi}{4}-2\pi=-\frac{3\pi}{4}.
\]
Por tanto,
\[
\Arg(-1-i)=-\frac{3\pi}{4}.
\]
Así,
\[
\boxed{
\operatorname{Log}(-1-i)
=
\frac12\ln2-\frac{3\pi}{4}i.
}
\]

\begin{ejercicio}
Determine todos los valores de $\log(-1-i)$ y señale el principal.
\end{ejercicio}

\textbf{Solución.}
A partir del argumento principal,
\[
\arg(-1-i)
=
-\frac{3\pi}{4}+2k\pi.
\]
Por tanto,
\[
\boxed{
\log(-1-i)
=
\frac12\ln2
+i\left(-\frac{3\pi}{4}+2k\pi\right),
\quad k\in\mathbb Z.
}
\]
El valor correspondiente a $k=0$ es
\[
\boxed{
\operatorname{Log}(-1-i)
=
\frac12\ln2-\frac{3\pi}{4}i.
}
\]

\begin{ejercicio}
Compare los valores principales de
\[
z_1=e^{i3\pi/4}
\qquad\text{y}\qquad
z_2=e^{i5\pi/4}.
\]
\end{ejercicio}

\textbf{Solución.}
Para $z_1$,
\[
\Arg z_1=\frac{3\pi}{4}.
\]
Como $|z_1|=1$,
\[
\Log z_1=\frac{3\pi}{4}i.
\]

Para $z_2$, el ángulo $5\pi/4$ no pertenece al intervalo principal. Restamos
$2\pi$:
\[
\frac{5\pi}{4}-2\pi=-\frac{3\pi}{4}.
\]
Por tanto,
\[
\Log z_2=-\frac{3\pi}{4}i.
\]

Así,
\[
\boxed{
\Log(e^{i3\pi/4})=\frac{3\pi}{4}i,
\qquad
\Log(e^{i5\pi/4})=-\frac{3\pi}{4}i.
}
\]

Este ejercicio permite observar que, para la rama principal,
\[
\Log(e^{i\theta})
\]
no es simplemente $i\theta$ para todo $\theta\in\mathbb R$: primero debemos
reducir $\theta$ al intervalo principal.

\begin{ejercicio}
Calcule
\[
\Log(e^{i7\pi/3}).
\]
\end{ejercicio}

\textbf{Solución.}
Reducimos el argumento módulo $2\pi$:
\[
\frac{7\pi}{3}-2\pi
=
\frac{\pi}{3}.
\]
Como $\pi/3\in(-\pi,\pi]$,
\[
\Arg(e^{i7\pi/3})=\frac{\pi}{3}.
\]
Además, el módulo es $1$, de modo que
\[
\boxed{
\Log(e^{i7\pi/3})=\frac{\pi}{3}i.
}
\]
Obsérvese que
\[
\boxed{
\Log(e^{i7\pi/3})\neq i\frac{7\pi}{3}.
}
\]

\begin{ejercicio}
Explique por qué no puede definirse una rama continua del logaritmo sobre
$\mathbb C\setminus\{0\}$.
\end{ejercicio}

\textbf{Solución.}
Una rama continua del logaritmo produciría, mediante su parte imaginaria, una
elección continua del argumento. Al recorrer una vez la circunferencia
\[
\gamma(t)=e^{it},\qquad0\leq t\leq2\pi,
\]
el argumento tendría que aumentar en $2\pi$. Sin embargo,
\[
\gamma(0)=\gamma(2\pi)=1.
\]
Una función univaluada tendría que asignar el mismo valor al punto inicial y
final, mientras que la continuidad del argumento a lo largo de la vuelta
exigiría valores separados por $2\pi$. Por ello tal rama no puede existir en
todo $\mathbb C\setminus\{0\}$.

\subsubsection{Resumen para la exposición}

La discusión puede sintetizarse en la siguiente cadena lógica:
\[
\boxed{
e^{z+2\pi i}=e^z
}
\]
implica
\[
\boxed{
\arg z=\theta+2k\pi
}
\]
y, por tanto,
\[
\boxed{
\log z=\ln|z|+i(\theta+2k\pi)
}
\]
es multivaluado.

Para obtener una función debemos seleccionar una rama:
\[
\boxed{
-\pi<\Arg z<\pi
}
\]
sobre
\[
\boxed{
\mathbb C\setminus(-\infty,0].
}
\]
Entonces
\[
\boxed{
\Log z=\ln|z|+i\Arg z
}
\]
es la rama principal analítica.

El corte no es un artificio puramente notacional: evita que podamos dar una
vuelta completa alrededor del origen dentro del dominio. De esta manera el
argumento puede escogerse de forma única y continua.

Finalmente, al acercarnos al corte por arriba y por abajo,
\[
\Arg z\to\pi
\qquad\text{y}\qquad
\Arg z\to-\pi,
\]
respectivamente, y el logaritmo presenta un salto de
\[
\boxed{2\pi i.}
\]


\subsection{Resolución de ecuaciones mediante el logaritmo complejo}

\begin{ejemplo}
Resolver
\[
e^z=3.
\]

Sea $z=x+iy$. Equivalentemente, podemos tomar el logaritmo complejo de $3$:
\[
z=\log3.
\]
Como $3=3e^{i2k\pi}$,
\[
\boxed{
z=\ln3+2k\pi i,\qquad k\in\mathbb Z.
}
\]

Comprobación:
\[
e^{\ln3+2k\pi i}
=3e^{2k\pi i}=3.
\]
\end{ejemplo}

\begin{ejemplo}
Resolver
\[
e^z=-4.
\]

Como
\[
-4=4e^{i(\pi+2k\pi)},
\]
tenemos
\[
\boxed{
z=\ln4+i(\pi+2k\pi),\qquad k\in\mathbb Z.
}
\]
\end{ejemplo}

\begin{ejemplo}
Resolver
\[
e^{2z}=1+i.
\]

\textbf{Paso 1. Tomar logaritmos complejos.}
\[
2z=\log(1+i).
\]
Ya sabemos que
\[
\log(1+i)
=
\frac12\ln2
+i\left(\frac{\pi}{4}+2k\pi\right).
\]

\textbf{Paso 2. Despejar $z$.}
\[
\boxed{
z=
\frac14\ln2
+i\left(\frac{\pi}{8}+k\pi\right),
\qquad k\in\mathbb Z.
}
\]
\end{ejemplo}

\begin{ejemplo}
Resolver
\[
e^{3z}=-8i.
\]

Como
\[
-8i=8e^{i(-\pi/2+2k\pi)},
\]
\[
3z=\ln8+i\left(-\frac{\pi}{2}+2k\pi\right).
\]
Dividiendo entre $3$:
\[
\boxed{
z=\frac{\ln8}{3}
+i\left(-\frac{\pi}{6}+\frac{2k\pi}{3}\right),
\qquad k\in\mathbb Z.
}
\]
Como $\ln8=3\ln2$:
\[
\boxed{
z=\ln2+i\left(-\frac{\pi}{6}+\frac{2k\pi}{3}\right).
}
\]
\end{ejemplo}

\subsection{Potencias complejas definidas mediante el logaritmo}

Una potencia con exponente complejo puede definirse mediante
\[
\boxed{
z^w=e^{w\log z}.
}
\]
Como $\log z$ es multivaluado, en general $z^w$ también puede ser
multivaluado.

\begin{ejemplo}
Calcular todos los valores de
\[
i^i.
\]

Sabemos que
\[
\log i=i\left(\frac{\pi}{2}+2k\pi\right).
\]
Entonces
\[
i^i=e^{i\log i}.
\]
Sustituyendo:
\[
i\log i
=
i\left[
i\left(\frac{\pi}{2}+2k\pi\right)
\right]
=
-\left(\frac{\pi}{2}+2k\pi\right).
\]
Por tanto,
\[
\boxed{
i^i=e^{-\pi/2-2k\pi},
\qquad k\in\mathbb Z.
}
\]

Un resultado notable es que \emph{todos estos valores son reales y positivos}.
El valor correspondiente al logaritmo principal es
\[
\boxed{
e^{-\pi/2}.
}
\]
\end{ejemplo}

\begin{ejemplo}
Calcular los valores de
\[
(-1)^i.
\]

Como
\[
\log(-1)=i(2k+1)\pi,
\]
entonces
\[
(-1)^i
=
e^{i\log(-1)}
=
e^{-(2k+1)\pi}.
\]
Así,
\[
\boxed{
(-1)^i=e^{-(2k+1)\pi},
\qquad k\in\mathbb Z.
}
\]
El valor principal es
\[
\boxed{e^{-\pi}.}
\]
\end{ejemplo}

\subsection{Ejercicios resueltos: logaritmo complejo}

\begin{ejercicio}
Calcular todos los valores de
\[
\log(-i).
\]
\end{ejercicio}

\textbf{Solución.}
\[
|-i|=1,\qquad
\operatorname{Arg}(-i)=-\frac{\pi}{2}.
\]
Por tanto,
\[
\boxed{
\log(-i)
=
i\left(-\frac{\pi}{2}+2k\pi\right),
\quad k\in\mathbb Z.
}
\]
El valor principal es
\[
\boxed{-\frac{\pi i}{2}.}
\]

\begin{ejercicio}
Calcular $\log(3+3i)$ y su valor principal.
\end{ejercicio}

\textbf{Solución.}
\[
|3+3i|=3\sqrt2,\qquad
\Arg(3+3i)=\frac{\pi}{4}.
\]
Entonces
\[
\boxed{
\log(3+3i)
=
\ln(3\sqrt2)
+i\left(\frac{\pi}{4}+2k\pi\right).
}
\]
Como
\[
\ln(3\sqrt2)=\ln3+\frac12\ln2,
\]
\[
\boxed{
\Log(3+3i)
=
\ln3+\frac12\ln2+\frac{\pi}{4}i.
}
\]

\begin{ejercicio}
Calcular $\log(-\sqrt3+i)$.
\end{ejercicio}

\textbf{Solución.}
El módulo es
\[
\sqrt{3+1}=2.
\]
El punto está en el segundo cuadrante y el ángulo de referencia es $\pi/6$:
\[
\Arg(-\sqrt3+i)=\frac{5\pi}{6}.
\]
Así,
\[
\boxed{
\log(-\sqrt3+i)
=
\ln2+i\left(\frac{5\pi}{6}+2k\pi\right),
\quad k\in\mathbb Z.
}
\]

\begin{ejercicio}
Calcular $\log(-2+2\sqrt3\,i)$.
\end{ejercicio}

\textbf{Solución.}
\[
|-2+2\sqrt3 i|
=\sqrt{4+12}=4.
\]
El punto está en el segundo cuadrante y
\[
\tan\theta=-\sqrt3,
\]
de donde
\[
\Arg z=\frac{2\pi}{3}.
\]
Por tanto,
\[
\boxed{
\log(-2+2\sqrt3 i)
=
\ln4+i\left(\frac{2\pi}{3}+2k\pi\right).
}
\]

\begin{ejercicio}
Resolver
\[
e^z=2i.
\]
\end{ejercicio}

\textbf{Solución.}
\[
2i=2e^{i(\pi/2+2k\pi)}.
\]
Por tanto,
\[
\boxed{
z=\ln2+i\left(\frac{\pi}{2}+2k\pi\right),
\qquad k\in\mathbb Z.
}
\]

\begin{ejercicio}
Resolver
\[
e^{2z}=-1.
\]
\end{ejercicio}

\textbf{Solución.}
\[
2z=i(\pi+2k\pi).
\]
Por tanto,
\[
\boxed{
z=i\left(\frac{\pi}{2}+k\pi\right),
\qquad k\in\mathbb Z.
}
\]

\begin{ejercicio}
Resolver
\[
e^{z-1}=1-i.
\]
\end{ejercicio}

\textbf{Solución.}
\[
1-i=\sqrt2\,e^{i(-\pi/4+2k\pi)}.
\]
Entonces
\[
z-1
=
\frac12\ln2
+i\left(-\frac{\pi}{4}+2k\pi\right),
\]
y
\[
\boxed{
z=
1+\frac12\ln2
+i\left(-\frac{\pi}{4}+2k\pi\right),
\quad k\in\mathbb Z.
}
\]

\begin{ejercicio}
Resolver
\[
e^{(1+i)z}=1.
\]
\end{ejercicio}

\textbf{Solución.}
Como
\[
\log1=2k\pi i,
\]
tenemos
\[
(1+i)z=2k\pi i.
\]
Por tanto,
\[
z=\frac{2k\pi i}{1+i}.
\]
Multiplicando por el conjugado:
\[
z=
\frac{2k\pi i(1-i)}{2}
=k\pi(i-i^2)
=k\pi(1+i).
\]
Así,
\[
\boxed{
z=k\pi(1+i),\qquad k\in\mathbb Z.
}
\]

\begin{ejercicio}
Calcular todos los valores de
\[
(1+i)^i.
\]
\end{ejercicio}

\textbf{Solución.}
Primero,
\[
\log(1+i)
=
\frac12\ln2
+i\left(\frac{\pi}{4}+2k\pi\right).
\]
Entonces
\[
(1+i)^i=e^{i\log(1+i)}.
\]
Multiplicando:
\[
i\log(1+i)
=
\frac{i}{2}\ln2
-\left(\frac{\pi}{4}+2k\pi\right).
\]
Por tanto,
\[
\boxed{
(1+i)^i
=
e^{-\pi/4-2k\pi}
e^{\,i(\ln2)/2},
\qquad k\in\mathbb Z.
}
\]
En forma trigonométrica:
\[
\boxed{
(1+i)^i
=
e^{-\pi/4-2k\pi}
\left[
\cos\left(\frac{\ln2}{2}\right)
+i\sin\left(\frac{\ln2}{2}\right)
\right].
}
\]

\subsection{Ejercicios adicionales para práctica}

Resolver detalladamente:

\begin{enumerate}
\item $\log(2)$.
\item $\log(-2)$.
\item $\log(2i)$.
\item $\log(1-i)$.
\item $\log(-1-i)$.
\item $\log(\sqrt3+i)$.
\item $\log(-3+3\sqrt3\,i)$.
\item $\log\left(\frac12+\frac{\sqrt3}{2}i\right)$.
\item $e^z=5$.
\item $e^z=-5i$.
\item $e^{2z}=4i$.
\item $e^{3z}=1$.
\item $e^{z+2}=1+i$.
\item $e^{(2-i)z}=1$.
\item Calcular todos los valores de $(-i)^i$.
\item Calcular todos los valores de $2^i$.
\item Calcular todos los valores de $(1-i)^i$.
\end{enumerate}

\subsection{Respuestas de los ejercicios adicionales}

\begin{enumerate}
\item
\[
\boxed{\log2=\ln2+2k\pi i.}
\]

\item
\[
\boxed{\log(-2)=\ln2+i(\pi+2k\pi).}
\]

\item
\[
\boxed{\log(2i)=\ln2+i\left(\frac{\pi}{2}+2k\pi\right).}
\]

\item
\[
\boxed{
\log(1-i)=\frac12\ln2
+i\left(-\frac{\pi}{4}+2k\pi\right).
}
\]

\item
\[
\boxed{
\log(-1-i)=\frac12\ln2
+i\left(-\frac{3\pi}{4}+2k\pi\right).
}
\]

\item
\[
\boxed{
\log(\sqrt3+i)
=\ln2+i\left(\frac{\pi}{6}+2k\pi\right).
}
\]

\item Como el módulo es $6$ y el argumento principal es $2\pi/3$,
\[
\boxed{
\log(-3+3\sqrt3 i)
=\ln6+i\left(\frac{2\pi}{3}+2k\pi\right).
}
\]

\item El número tiene módulo $1$ y argumento $\pi/3$:
\[
\boxed{
\log\left(\frac12+\frac{\sqrt3}{2}i\right)
=i\left(\frac{\pi}{3}+2k\pi\right).
}
\]

\item
\[
\boxed{z=\ln5+2k\pi i.}
\]

\item
\[
\boxed{
z=\ln5+i\left(-\frac{\pi}{2}+2k\pi\right).
}
\]

\item
\[
2z=\ln4+i\left(\frac{\pi}{2}+2k\pi\right),
\]
por tanto
\[
\boxed{
z=\ln2+i\left(\frac{\pi}{4}+k\pi\right).
}
\]

\item
\[
3z=2k\pi i,
\]
por lo que
\[
\boxed{z=\frac{2k\pi}{3}i.}
\]

\item
\[
z+2=\frac12\ln2+i\left(\frac{\pi}{4}+2k\pi\right),
\]
de modo que
\[
\boxed{
z=-2+\frac12\ln2
+i\left(\frac{\pi}{4}+2k\pi\right).
}
\]

\item
\[
(2-i)z=2k\pi i,
\]
de donde
\[
z=\frac{2k\pi i}{2-i}
=\frac{2k\pi i(2+i)}5
=\boxed{\frac{2k\pi}{5}(-1+2i)}.
\]

\item
\[
\log(-i)=i\left(-\frac{\pi}{2}+2k\pi\right),
\]
por tanto
\[
\boxed{
(-i)^i=e^{\pi/2-2k\pi}.
}
\]

\item
\[
\log2=\ln2+2k\pi i.
\]
Entonces
\[
2^i=e^{i\ln2-2k\pi},
\]
es decir,
\[
\boxed{
2^i=e^{-2k\pi}
\left(\cos(\ln2)+i\sin(\ln2)\right).
}
\]

\item
\[
\log(1-i)
=
\frac12\ln2+i\left(-\frac{\pi}{4}+2k\pi\right).
\]
Multiplicando por $i$:
\[
i\log(1-i)
=
\frac{i}{2}\ln2+\frac{\pi}{4}-2k\pi.
\]
Por tanto,
\[
\boxed{
(1-i)^i
=
e^{\pi/4-2k\pi}
e^{i(\ln2)/2}.
}
\]
\end{enumerate}

\subsection{Síntesis conceptual}

Para $z\neq0$, las ideas esenciales pueden resumirse en:
\[
\boxed{
\log z=\ln|z|+i(\arg z+2k\pi),\qquad k\in\mathbb Z.
}
\]

El logaritmo principal selecciona un argumento:
\[
\boxed{
\Log z=\ln|z|+i\Arg z.
}
\]

La multivaluación no es una anomalía algebraica, sino una consecuencia directa
de la periodicidad
\[
e^{z+2\pi i}=e^z.
\]

Finalmente, el logaritmo complejo permite resolver ecuaciones exponenciales y
definir potencias complejas:
\[
e^w=z
\quad\Longleftrightarrow\quad
w\in\log z,
\]
y
\[
z^w=e^{w\log z}.
\]

Estos hechos constituyen el puente natural hacia el estudio posterior de
ramas, funciones analíticas, derivación compleja y funciones elementales en
variable compleja.


\section{Cierre}

Los ejercicios anteriores permiten separar tres cuestiones fundamentales:

\begin{enumerate}
\item \textbf{Existencia.} El TFA garantiza que todo polinomio no constante
con coeficientes complejos posee una raíz en $\C$.

\item \textbf{Factorización.} Combinado con el teorema del factor, permite
escribir un polinomio de grado $n$ como producto de $n$ factores lineales,
contando multiplicidades.

\item \textbf{Cálculo explícito.} El TFA no proporciona por sí mismo un
algoritmo para calcular las raíces. En los problemas anteriores las soluciones
se obtienen gracias a estructuras particulares: forma polar, raíces
$n$-ésimas, sustituciones como $u=z^m$, factorización y simetría por
conjugación.
\end{enumerate}

Por ello es importante que el estudiante distinga entre afirmar
\emph{que las raíces existen} y disponer de un procedimiento efectivo para
\emph{calcularlas}.

\end{document}
